% part: intuitionistic-logic
% chapter: introduction
% section: natural-deduction

\documentclass[../../../include/open-logic-section]{subfiles}

\begin{document}

\olfileid{int}{int}{ntd}

\olsection{自然演绎}

不含 $\FalseCl$ 规则的自然演绎是直觉主义逻辑的标准推导系统。
我们在此复述这些规则，并利用 BHK 解释来阐明其动机。
在每种情况下，我们都可以这样来理解一条规则：它允许我们得出，如果前提有构造，那么结论也有构造。

由于自然演绎推导可以有未解除的假设，我们应当将这样的推导——例如，从未解除的假设 $\Gamma$ 推导 $!A$——视为一个函数，它将所有 $!B \in \Gamma$ 的构造转化为 $!A$ 的一个构造。
如果存在一个没有任何未解除假设的 $!A$ 的推导，那么根据 BHK 解释，$!A$ 就有一个构造。
不过，为了讨论方便，我们在不需要时会略去 $\Gamma$。

一个假设 $!A$ 本身就是从其未解除假设 $!A$ 推导 $!A$ 的推导。
这与 BHK 解释一致：构造上的恒等函数将 $!A$ 的任何构造转化为 $!A$ 的一个构造。

\subsection{合取}

\begin{defish}
\AxiomC{$!A$}
\AxiomC{$!B$}
\RightLabel{\Intro{\land}}
\BinaryInfC{$!A \land !B$}
\DisplayProof
\hfill
\begin{tabular}{l}
\AxiomC{$!A \land !B$}
\RightLabel{\Elim{\land}}
\UnaryInfC{$!A$}
\DisplayProof
\\[3ex]
\AxiomC{$!A \land !B$}
\RightLabel{\Elim{\land}}
\UnaryInfC{$!B$}
\DisplayProof
\end{tabular}
\end{defish}

假设我们分别有 $!A_1$ 和 $!A_2$ 的构造 $N_1$ 和 $N_2$。
那么我们也有 $!A_1 \land !A_2$ 的一个构造，即有序对 $\tuple{N_1, N_2}$。

在 BHK 解释下，$!A_1 \land !A_2$ 的一个构造是一个有序对 $\tuple{N_1, N_2}$。
因此，假设我们有了这样一个有序对，那么我们也就有了每个合取支的一个构造：
$N_1$ 是 $!A_1$ 的一个构造，而 $N_2$ 是 $!A_2$ 的一个构造。

\subsection{条件句}

\begin{defish}
  \AxiomC{$\Discharge{!A}{u}$}
  \DeduceC{$!B$}
  \DischargeRule{$\Intro{\lif}$}{u}
  \UnaryInfC{$!A \lif !B$}
  \DisplayProof
\hfill
  \AxiomC{$!A \lif !B$}
  \AxiomC{$!A$}
  \RightLabel{\Elim{\lif}}
  \BinaryInfC{$!B$}
  \DisplayProof
\end{defish}

如果我们有一个从未解除假设 $!A$ 推导出 $!B$ 的推导，那么就存在一个函数 $f$，它将 $!A$ 的构造转化为 $!B$ 的构造。
这个函数本身也就是 $!A \lif !B$ 的一个构造。
所以，如果 $\Intro\lif$ 的前提在给定 $!A$ 的一个构造的条件下有构造，那么结论 $!A \lif !B$ 就有一个构造。

另一方面，假设我们有 $!A$ 的一个构造 $N$ 以及 $!A \lif !B$ 的一个构造 $f$。
$!A \lif !B$ 的一个构造是一个将 $!A$ 的构造转化为 $!B$ 的构造的函数。
因此，$f(N)$ 就是 $!B$ 的一个构造，也就是说，$\Elim\lif$ 的结论有一个构造。

\subsection{析取}

\begin{defish}
\begin{tabular}{l}
\AxiomC{$!A$}
\RightLabel{\Intro{\lor}}
\UnaryInfC{$!A \lor !B$}
\DisplayProof
\\[3ex]
\AxiomC{$!B$}
\RightLabel{\Intro{\lor}}
\UnaryInfC{$!A \lor !B$}
\DisplayProof
\end{tabular}
\hfill
\AxiomC{$!A \lor !B$}
\AxiomC{$\Discharge{!A}{n}$}
\DeduceC{$!C$}
\AxiomC{$\Discharge{!B}{n}$}
\DeduceC{$!C$}
\DischargeRule{\Elim{\lor}}{n}
\TrinaryInfC{$!C$}
\DisplayProof
\end{defish}

如果我们有 $!A_i$ 的一个构造 $N_i$，我们可以将其转化为 $!A_1 \lor !A_2$ 的一个构造 $\tuple{i, N_i}$。
另一方面，假设我们有 $!A_1 \lor !A_2$ 的一个构造，即一个有序对 $\tuple{i, N_i}$，其中 $N_i$ 是 $!A_i$ 的一个构造；
以及函数 $f_1$ 和 $f_2$，它们分别将 $!A_1$ 和 $!A_2$ 的构造转化为 $!C$ 的构造。
那么 $f_i(N_i)$ 就是 $!C$ 的一个构造，也就是 $\Elim\lor$ 结论的构造。

\subsection{荒谬}

\begin{defish}
  \AxiomC{$\lfalse$}
  \RightLabel{\FalseInt}
  \UnaryInfC{$!A$}
  \DisplayProof
\end{defish}

如果我们有一个从一组未解除假设 $!B_1$, \dots, $!B_n$ 推导出 $\lfalse$ 的推导，那么就存在一个函数 $f(M_1,
\dots, M_n)$，它将 $!B_1$, \dots, $!B_n$ 的构造转化为 $\lfalse$ 的构造。
由于 $\lfalse$ 没有构造，$!B_1$, \dots, $!B_n$ 也不可能同时都有构造。
因此，函数 $f$ 也就具有如下性质：\emph{如果} $M_1$, \dots, $M_n$ 分别是 $!B_1$, \dots, $!B_n$ 的构造，\emph{那么} $f(M_1, \dots, M_n)$ 就是 $!A$ 的一个构造。

\subsection{关于 $\lnot$ 的规则}

由于 $\lnot !A$ 被定义为 $!A \lif \lfalse$，严格来说我们不需要关于 $\lnot$ 的专门规则。
但如果我们确实需要，它们的形式如下：

\begin{defish}
\AxiomC{$\Discharge{!A}{n}$}
\noLine
\DeduceC{$\lfalse$}
\DischargeRule{\Intro{\lnot}}{n}
\UnaryInfC{$\lnot !A$}
\DisplayProof
\hfill
\AxiomC{$\lnot !A$}
\AxiomC{$!A$}
\RightLabel{\Elim{\lnot}}
\BinaryInfC{$\lfalse$}
\DisplayProof
\end{defish}


\subsection{推导示例}

\begin{enumerate}
\item $\Proves !A \lif (\lnot !A \lif \lfalse)$，即 $\Proves !A
  \lif ((!A \lif \lfalse) \lif \lfalse)$
  \begin{prooftree}
    \AxiomC{$\Discharge{!A}{2}$}
    \AxiomC{$\Discharge{!A \lif \lfalse}{1}$}
    \RightLabel{\Elim\lif}
    \BinaryInfC{$\lfalse$}
    \DischargeRule{\Intro\lif}{1}
    \UnaryInfC{$(!A \lif \lfalse)\lif \lfalse$}
    \DischargeRule{\Intro\lif}{2}
    \UnaryInfC{$!A \lif (!A \lif \lfalse) \lif \lfalse$}
  \end{prooftree}

\item $\Proves ((!A \land !B) \lif !C) \lif (!A \lif (!B \lif !C))$
  \begin{prooftree}
    \AxiomC{$\Discharge{(!A \land !B) \lif !C}{3}$}
    \AxiomC{$\Discharge{!A}{2}$}
    \AxiomC{$\Discharge{!B}{1}$}
    \RightLabel{\Intro\land}
    \BinaryInfC{$!A \land !B$}
    \RightLabel{\Elim\lif}
    \BinaryInfC{$!C$}
    \DischargeRule{\Intro\lif}{1}
    \UnaryInfC{$!B \lif !C$}
    \DischargeRule{\Intro\lif}{2}
    \UnaryInfC{$!A \lif (!B \lif !C)$}
    \DischargeRule{\Intro\lif}{3}
    \UnaryInfC{$((!A \land !B) \lif !C) \lif (!A \lif (!B \lif !C))$}
  \end{prooftree}

\item $\Proves \lnot(!A \land \lnot !A)$，即 $\Proves (!A \land (!A
  \lif \lfalse))\lif \lfalse$
  \begin{prooftree}
    \AxiomC{$\Discharge{!A \land (!A \lif \lfalse)}{1}$}
    \RightLabel{\Elim\land}
    \UnaryInfC{$!A \lif \lfalse$}
    \AxiomC{$\Discharge{!A \land (!A \lif \lfalse)}{1}$}
    \RightLabel{\Elim\land}
    \UnaryInfC{$!A$}
    \RightLabel{\Elim\lif}
    \BinaryInfC{$\lfalse$}
    \DischargeRule{\Intro\lif}{1}
    \UnaryInfC{$(!A \land (!A \lif \lfalse))\lif \lfalse$}
  \end{prooftree}

\item $\Proves \lnot\lnot(!A \lor \lnot !A)$，即 $\Proves ((!A \lor
  (!A \lif \lfalse))\lif \lfalse)\lif \lfalse$
  \begin{prooftree}\hskip-3em
    \AxiomC{$\Discharge{(!A \lor (!A \lif \lfalse))\lif \lfalse}{2}$}
    \AxiomC{$\Discharge{(!A \lor (!A \lif \lfalse))\lif \lfalse}{2}$}
    \AxiomC{$\Discharge{!A}{1}$}
    \RightLabel{\Intro\lor}
    \UnaryInfC{$!A \lor (!A \lif \lfalse)$}
    \RightLabel{\Elim\lif}
    \BinaryInfC{$\lfalse$}
    \DischargeRule{\Intro\lif}{1}
    \UnaryInfC{$!A \lif \lfalse$}
    \RightLabel{\Intro\lor}
    \UnaryInfC{$!A \lor (!A \lif \lfalse)$}
    \RightLabel{\Elim\lif}
    \insertBetweenHyps{\hskip -4em}
    \BinaryInfC{$\lfalse$}
    \DischargeRule{\Intro\lif}{2}
    \UnaryInfC{$((!A \lor (!A \lif \lfalse))\lif \lfalse)\lif \lfalse$}
  \end{prooftree}
\end{enumerate}

\begin{prop}
  若在直觉主义逻辑的自然演绎中有 $\Gamma \Proves !A$，则在经典逻辑中也有 $\Gamma \Proves !A$。
  特别地，如果 $!A$ 是一个直觉主义定理，那么它也是一个经典定理。
\end{prop}

\begin{proof}
  直觉主义自然演绎的每一条规则也是经典自然演绎的规则，因此直觉主义逻辑中的每一个推导本身也是经典逻辑中的一个推导。
\end{proof}

\begin{prob}
  给出下列公式在直觉主义逻辑中的推导：
  \begin{enumerate}
    \item $(\lnot !A \lor !B) \lif (!A \lif !B)$
    \item $\lnot\lnot\lnot !A \lif \lnot !A$
    \item $\lnot\lnot (!A \land !B) \liff (\lnot\lnot !A\land \lnot \lnot !B)$
    \item $\lnot (!A \lor !B) \liff (\lnot !A \land \lnot !B)$
    \item $(\lnot !A \lor \lnot !B) \lif \lnot (!A \land !B)$
    \item $\lnot\lnot ( !A \land !B) \lif  (\lnot\lnot !A \lor
    \lnot\lnot !B)$
  \end{enumerate}
\end{prob}

\end{document}